\documentclass{article}

\usepackage{amsmath,amssymb}
\usepackage{xurl}
\usepackage[hidelinks]{hyperref}
\setlength{\emergencystretch}{2em}

% Shared commands for notes in docs/paper-gaps/.

\DeclareUnicodeCharacter{2080}{\ifmmode{}_{0}\else\textsubscript{0}\fi}
\DeclareUnicodeCharacter{2081}{\ifmmode{}_{1}\else\textsubscript{1}\fi}
\DeclareUnicodeCharacter{2082}{\ifmmode{}_{2}\else\textsubscript{2}\fi}
\DeclareUnicodeCharacter{2099}{\ifmmode{}_{n}\else\textsubscript{n}\fi}

\newcommand{\Lean}{\textsc{Lean}}
\ExplSyntaxOn
\NewDocumentCommand{\leanid}{m}{
  \ifmmode
    \textsf{\detokenize{#1}}
  \else
    \group_begin:
    \urlstyle{sf}
    \tl_set:Nn \l_tmpa_tl {#1}
    \tl_replace_all:Nnn \l_tmpa_tl {\_} {_}
    \exp_args:NV \path \l_tmpa_tl
    \group_end:
  \fi
}
\ExplSyntaxOff
\providecommand{\tr}{\operatorname{tr}}
\newcommand{\tnleanIssueBase}{https://github.com/LionSR/TNLean/issues/}
\newcommand{\tnleanPrBase}{https://github.com/LionSR/TNLean/pull/}
\newcommand{\tnleanDocBase}{https://sirui-lu.com/TNLean/docs/}
\newcommand{\ghissue}[1]{\href{\tnleanIssueBase#1}{GitHub issue \##1}}
\newcommand{\ghpr}[1]{\href{\tnleanPrBase#1}{PR \##1}}
\newcommand{\doclink}[1]{\href{\tnleanDocBase#1.html}{\path{#1}}}

% Dirac notation and the identity operator, as in blueprint/src/macros/common.tex.
% \providecommand keeps note-local definitions authoritative.
\usepackage{dsfont}
\providecommand{\ket}[1]{|#1\rangle}
\providecommand{\bra}[1]{\langle #1|}
\providecommand{\braket}[2]{\langle #1 | #2 \rangle}
\providecommand{\ketbra}[2]{|#1\rangle\!\langle #2|}
\providecommand{\Id}{\mathds{1}}
\providecommand{\one}{\mathds{1}}
\providecommand{\C}{\mathbb{C}}
\providecommand{\MN}[1]{M_{#1}(\C)}
\providecommand{\MD}{M_D(\C)}

% Verdict marker: \gapnote{<kind>}{<status>}. Kind is the policy
% classification (clarification, local-correction, scope-restriction,
% unfaithful, false-source, open-gap); status is open, wip (elimination
% underway), resolved, or historical. Machine-read by the site index;
% prints a verdict line.
\newcommand{\gapnote}[2]{\par\noindent\textbf{Verdict:} #1 (#2).\par}


% Fallbacks keep this note compilable when a different command.tex is found first.
\providecommand{\leanid}[1]{\textsf{\detokenize{#1}}}
\providecommand{\gapnote}[2]{\par\noindent\textbf{Verdict:} #1 (#2).\par}
\providecommand{\doclink}[1]{\path{#1}}

\title{The Inverse-Gap Logarithm in the Area Law of Arad, Kitaev, Landau and Vazirani}
\author{}
\date{2026-11-15}

\begin{document}
\maketitle
\gapnote{local-correction}{resolved}

\section{The source statement}

Arad, Kitaev, Landau and Vazirani~\cite{Arad2013AreaLaw},
arXiv:1301.1162v1, \path{1dfrustrated34.tex}, state the following.
\begin{itemize}
    \item Theorem~4.3 (label \texttt{t:ent}, lines~820--846): for a
          frustration-free nearest-neighbour chain of local dimension $d$ with
          terms $0\le H_i\le1$, a unique ground state $\Gamma$ and gap
          $\varepsilon$, the entanglement entropy of $\Gamma$ across every cut
          is $O(\log^3d/\varepsilon)$.
    \item The general theorem of Section~6 (lines~995--1000): the same bound
          $O(\log^3d/\varepsilon)$ for frustrated chains with gap
          $\varepsilon\le1$.
    \item Lemma~6.3 (label \texttt{l:5}, lines~1003--1052): a sequence of
          states $\psi_i$ with Schmidt ranks $R_i$, where
          $\log R_0=O(\log^3d/\varepsilon)$.
\end{itemize}
The abstract (line~162) and both theorems state $O(\log^3d/\varepsilon)$; the
introduction (line~214) states $\tilde O(\log^3d/\varepsilon)$.

All three statements rest on one approximate ground-state projector. Clip the
two exterior blocks of the Hamiltonian at height $t=O(1/\varepsilon)$, so that
the clipped Hamiltonian, with $s$ terms in the window around the cut, has norm
at most $w=s+2t$ and gap $\gamma$ of order $\varepsilon$. Let $K$ be the
Chebyshev polynomial of degree $\ell$ of the clipped Hamiltonian, normalized at
the ground energy. Lemma~4.1 (lines~530--583) gives the squared shrinking
factor
\begin{align}
    \Delta\le4\exp\bigl(-4\ell\sqrt{\gamma/w}\bigr),
    \label{eq:aklv13_igl_shrink}
\end{align}
and Lemma~4.2 (line~590) bounds the operator Schmidt rank of $K$ across the
cut by
\begin{align}
    D\le\bigl(d(\ell+1)\bigr)^{c(\lfloor\ell/s\rfloor+\lfloor\sqrt\ell\rfloor+1)},
    \label{eq:aklv13_igl_rank}
\end{align}
with an absolute constant $c$; the count with $\ell+1$ is the corrected one of the QICLean note
\url{https://sirui-lu.com/QICLean/paper-gaps/aklv13_interpolation_and_word_count.pdf},
and the proved constant is
$c=16$, or $c=32$ at the shifted cut of the frustration-free filter. The
entropy is bounded through Corollary~2.4 of Ref.~\cite{Arad2013AreaLaw}:
if $D\Delta\le1/2$ then $S\le O(\log D)$. The frustrated argument starts its
sequence from such a pair, with $R_0=D$.

\section{Where the inverse-gap logarithm enters}

The proof of Theorem~4.3 takes $\ell=s^2/2$ and $s=O(\log^2d/\varepsilon)$.
With $t$ of order $1/\varepsilon$ and $s\ge1/\varepsilon$, one has
$w=O(s)$ and $-\log\Delta$ is of order $s^{3/2}\varepsilon^{1/2}$. The rank
bound~(\ref{eq:aklv13_igl_rank}) at $\ell$ of order $s^2$ gives
\begin{align}
    \log D=O\bigl(s(\log d+\log s)\bigr).
    \label{eq:aklv13_igl_square_degree_rank}
\end{align}
The source keeps this logarithm: lines~840--841 display the factor
$\log(s^2+s)$, and lines~842--843 drop it when $s$ is chosen;
lines~831--836 quote $\Delta$ and $D$ from a commented earlier form of
Lemma~4.2 (lines~700--707), with the same logarithm. Since
$s\ge1/\varepsilon$, the dropped term $\log s$
of~(\ref{eq:aklv13_igl_square_degree_rank}) is at least
$\log(1/\varepsilon)$. The condition $D\Delta\le1/2$ then requires
$\sqrt{s\varepsilon}\ge C(\log d+\log s)$, so $s$ must be of order
$(\log d+\log(1/\varepsilon))^2/\varepsilon$ rather than
$\log^2d/\varepsilon$. With this choice,
\begin{align}
    \log D=O\Bigl(\frac{(\log d+\log(1/\varepsilon))^3}{\varepsilon}\Bigr).
    \label{eq:aklv13_igl_source_rate}
\end{align}
The rate~(\ref{eq:aklv13_igl_source_rate}) is the rate proved in the
repository; see Section~\ref{sec:igl_formal}.

\section{No parameter choice avoids it}

The inverse-gap logarithm is not an artefact of the particular choice
$\ell=s^2/2$. Fix $c\ge4$, $q=d\ge2$, natural numbers $\ell$ and $s\ge1$,
and reals $0\le\gamma\le g$ and $w\ge s$. Write $D$ and $\Delta$ for the right
sides of (\ref{eq:aklv13_igl_rank}) and~(\ref{eq:aklv13_igl_shrink}).
\begin{quote}
    If $D\Delta\le1$, then
    \begin{align}
        \log D>\frac{c^3}{16g}\Bigl(\log q+2\log\frac1g\Bigr)^3.
        \label{eq:aklv13_igl_lower}
    \end{align}
\end{quote}
\emph{Proof.} Put $N=\lfloor\ell/s\rfloor+\lfloor\sqrt\ell\rfloor+1$,
$\Lambda=\log(q(\ell+1))$ and $X=4\ell\sqrt{\gamma/w}$. If $\ell=0$ then
$D\Delta=4D\ge4$, so $\ell\ge1$ and $\Lambda\ge\log4>1$. The hypothesis gives
$cN\Lambda=\log D<X$, and $X^2\le16\ell^2g/s$. The floors give $\ell<sN$ and
$\ell<N^2$. Squaring and using $\ell<sN$ gives
$c^2N\Lambda^2<16g\ell$; using $\ell<N^2$ then gives
$c^2\Lambda^2<16gN$. Hence $\log D=cN\Lambda>c^3\Lambda^3/(16g)$. Combining
the two inequalities gives $\ell>(c\Lambda)^4/(256g^2)\ge1/g^2$, so
$\Lambda=\log q+\log(\ell+1)\ge\log q+2\log(1/g)$, and the cube is monotone.
\hfill$\square$

The condition $\gamma\le g$ is the form in which the gap enters every
application: in the frustration-free case the clipped Hamiltonian $K'$
satisfies $0\le K'\le H$ with the same ground vector, so its gap is at most
$g$ by the min-max principle; elsewhere the source uses a gap of order
$\varepsilon$, and a bound $\gamma\le\kappa g$ reduces to the statement at gap
$\kappa g$, changing only the constant. The condition $w\ge s$ holds for the
ceiling $w=s+2t$ of the source.

The same conclusion holds for the commented form $(d\ell)^{O(\ell/s+s)}$ of the
rank estimate: the term $\ell/s$ forces $s\ge c'\Lambda^2/g$, so
$\log D\ge c''s\Lambda\ge c'''\Lambda^3/g$, and $\ell\ge c''''/g^2$ again.

Consequently, for fixed $q$ the forced cost is of order
$\log^3(1/g)/g$, which is not bounded by any multiple of $\log^3q/g$: for every
$C$ there is $0<g\le1$ such that every favorable pair at gap $g$ has
$\log D>C\log^3q/g$. Since the source bounds the entropy by a multiple of
$\log D$ for such a pair, its argument cannot give $O(\log^3d/\varepsilon)$
with an absolute constant. The two certificates are close to optimal for the
method. The Chebyshev polynomial is the minimax polynomial normalized at the
ground energy, so (\ref{eq:aklv13_igl_shrink}) can be improved only in the
constant of the exponent. The binomial count behind (\ref{eq:aklv13_igl_rank})
is at least $s^s$ at $\ell=s^2$, as recorded in
\path{docs/paper-gaps/aklv13_frustration_free_mps_rate.tex}.
Removing the inverse-gap logarithm would therefore need a different rank
mechanism, not a different choice of parameters.

These are statements about the two upper-bound certificates. They give no
lower bound on the actual operator Schmidt rank of $K$ and do not refute the
printed entropy bound, which may hold by another argument.

\section{What is formalized}
\label{sec:igl_formal}

The inequality~(\ref{eq:aklv13_igl_lower}) is
\leanid{TNLean.MPS.AreaLaw.aklv_certificate_log_rank_gt}, and the
non-uniformity in $g$ is
\leanid{TNLean.MPS.AreaLaw.aklv_certificate_not_uniform_log_cube}, both in
\doclink{TNLean/MPS/AreaLaw/ChebyshevAGSP/CertificateRateLowerBound}.

The entropy bounds proved from the source's construction are, with
$u=1+\log q+\log(1/g)$,
\begin{align}
    S_k(\Omega)\le\frac Cg u^3
    \label{eq:aklv13_igl_formal_entropy}
\end{align}
at every cut, with a universal constant $C$: for frustration-free chains in
\leanid{TNLean.MPS.AreaLaw.exists_aklv_frustrationFree_physical_entropy_bound}
(\doclink{TNLean/MPS/AreaLaw/ChebyshevAGSP/FrustrationFree/FrustrationFreeEntropy}), and for frustrated
chains in \leanid{TNLean.MPS.AreaLaw.exists_aklv_physical_entropy_bound}
(\doclink{TNLean/MPS/AreaLaw/ChebyshevAGSP/PhysicalEntropy}). Since $\log q\ge\log2$,
$u\le(1+1/\log2)(\log q+\log(1/g))$, so the bounds~(\ref{eq:aklv13_igl_formal_entropy}) match
(\ref{eq:aklv13_igl_lower}) up to the universal constant. The lemma
\leanid{TNLean.MPS.AreaLaw.aklv_gap_logarithm_not_uniformly_absorbable}
records that $u^3$ is not a bounded multiple of $\log^32$, and
\leanid{TNLean.MPS.AreaLaw.aklv_cubic_cost_of_fixed_gap} gives the printed rate
with a constant depending on a lower bound $g_0\le g$.

\section{Verdict}

A local correction, resolved. Within the source's argument the entropy bound is
$O((\log d+\log(1/\varepsilon))^3/\varepsilon)$, and this is formalized. The
printed $O(\log^3d/\varepsilon)$ of Theorem~4.3, of the general theorem of
Section~6 and of Lemma~6.3 is not formalized: it drops the factor $\log(s^2+s)$
at lines~840--843, and no parameter choice in the source's certificates
recovers it. The tilde in the introduction (line~214) covers the factor.

\bibliographystyle{plain}
\bibliography{references}

\end{document}
